\section{Justification and explanation of the models}
\label{sec:justification}

Each subsection says what the model is asked, why the technique fits, what
had to be done to the data, and which of the models built in the notebook
was kept and why. The numbers are the ones printed by the executed
notebooks in the repository.

\subsection{Bitcoin price: ARIMA}

\paragraph{The question.} What will bitcoin cost in the next few days?

\paragraph{Why ARIMA.} There are no predictors here, only the history of
the price, and ARIMA is the technique of the course for forecasting a
series from its own past.

\paragraph{What the data needed.} The price grew twenty-fold over the four
years, and its daily jumps grew with it: a 10\,\% move was 20 dollars in 2015
and 250 in 2017. ARIMA needs a series with a stable mean and a stable
spread, so two steps were taken. Taking the logarithm makes the spread
stable, because it turns dollar changes into percentage changes. Taking
the difference between one day and the next removes the trend. After those
two steps the statistical test for stationarity passes.

\paragraph{What the model found.} The daily percentage change of bitcoin
has almost no memory: knowing today's change says almost nothing about
tomorrow's. That is what economists call a random walk, and it means the
best guess for tomorrow is roughly today's price. The notebook tried
sixteen ARIMA orders and kept the one with the best AIC score,
ARIMA(3,1,3), but none of its coefficients is statistically different from
zero, so it is really a random walk with a few small extra terms.

\paragraph{What is kept.} The ARIMA is kept because, unlike a plain
``same as today'' answer, it gives a range around the forecast that widens
with the horizon. The assistant reads out the number and the range, and
Section~\ref{sec:results} shows honestly that the number itself is no
better than repeating the last price.

\subsection{Used car price: multiple linear regression}

\paragraph{The question.} How much does a used car sell for, given its
showroom price, year, mileage, fuel, seller and transmission?

\paragraph{Why linear regression.} A dealer wants to know how much each
fact is worth, for example how much a year of age costs. A linear
regression gives exactly that: one coefficient per fact, in units of
price.

\paragraph{What the data needed.} Most cars are cheap and a few are very
expensive, so prices are lopsided. Fitting the price directly, the model
gets the expensive cars badly wrong and its errors grow with the price.
Fitting the logarithm of the price instead, with the logarithm of the
showroom price and of the mileage, fixes that; the coefficients then read
as percentages, which is the natural way to talk about depreciation.

\paragraph{What is kept.} The logarithmic model, with an $R^2$ of 0.954
against 0.807 for the direct model and half the error. Cars outside the
range seen in training (a much more expensive or much older car) are not
predicted, because a straight line knows nothing beyond its data. The
coefficient for CNG rests on two cars and is reported as unknown.

\subsection{Telecom churn: neural network}

\paragraph{The question.} Will this customer cancel the service?

\paragraph{Why a neural network.} A multilayer perceptron can combine the
30 facts about a customer (contract, services, bill, tenure) in ways a
straight line cannot. The network has two hidden layers of 32 and 16
neurons and one output neuron that gives a probability.

\paragraph{Why ReLU.} The activation of the hidden neurons is set
explicitly. The notebook tries the four that scikit-learn offers
(identity, logistic, tanh and ReLU) with the same architecture, on the
training data alone, and they land within one standard deviation of each
other: a cross-validated AUC between 0.830 and 0.839 with a fold moving it
by about 0.014. The identity row matches the rest, which says the churn
signal is largely linear, but a network with no non-linearity is a logistic
regression and is not kept. Among the non-linear three ReLU is kept because
its gradient does not fade on the positive side, and the other two buy
nothing measurable here in exchange for their saturating sums.

\paragraph{What the data needed.} Only about one customer in four leaves,
so the network learns to say ``stays'' most of the time: at the usual cut
of 0.5 it reaches 79\,\% accuracy but catches only half of the customers
who leave. The notebook sweeps the threshold and shows that the ranking
behind the probabilities is good (AUC 0.84): lowering the cut to 0.3 finds
three leavers out of four at the price of twice the false alarms.

\paragraph{What is kept.} The network at the 0.5 threshold, which is the
most accurate cut, with the probability stored so the company can lower the
threshold when a retention call is cheap next to a lost customer. The
weak number is stated plainly: about four customers in ten on the list
would not have left, so the list is a set of customers worth a call, not a
verdict.

\subsection{Wine quality: $k$ nearest neighbours}

\paragraph{The question.} How good is a wine, from its laboratory sheet?

\paragraph{Why $k$-NN.} It needs no equation: to score a new wine it looks
up the $k$ most similar wines in the training set and takes the most common
score among them.

\paragraph{What the data needed.} ``Similar'' means close in every
measurement, but the measurements use very different scales (sulphur
dioxide in the hundreds, density around 1). Without putting them on the
same scale the distance is decided by sulphur dioxide alone, so every
measurement was standardised. The second problem is the target: scores go
from 3 to 9 but almost every wine is a 5, 6 or 7, and a vote among
neighbours never picks the rare scores. Grouping the scores into low,
medium and high gives every class enough wines to win a vote.

\paragraph{What is kept.} The three-band model with $k = 39$. It gets the
band right 59\,\% of the time, against 44\,\% for always answering
``medium'', and its mistakes are between neighbouring bands. The honest
reading is that the laboratory sheet explains only part of what tasters
score, so the assistant gives a band, not a number.

\subsection{Flight delay: decision tree}

\paragraph{The question.} How late will this flight arrive?

\paragraph{Why a decision tree.} A tree asks a chain of yes/no questions
(``does it depart after 5 pm?'', ``is it this airline?'') and predicts the
average delay of the flights that answered the same way. Its rules can be
read, so the assistant could explain why it expects a delay.

\paragraph{What the data needed.} The file contains the departure delay
of every flight. A tree that reads it predicts the arrival delay almost
perfectly ($R^2$ of 0.95), but that is cheating: if you already know the
plane left forty minutes late you do not need a model. The honest tree only
reads what is known before leaving home: airline, route, month, weekday,
scheduled hour and distance.

\paragraph{What is kept.} The honest tree, with an $R^2$ of 0.05 and an
average error of 36 minutes. The notebook says plainly that this is a weak
model, and explains why: delays come from weather, breakdowns, late
inbound aircraft and airport congestion, none of which is in the schedule.
The assistant gives the expected delay together with that error.

\subsection{Stroke risk: logistic regression}

\paragraph{The question.} Is this person at risk of a stroke?

\paragraph{Why logistic regression.} It turns a set of facts (age, blood
pressure, heart disease, glucose, body mass, smoking) into a probability
between 0 and 1, and a threshold on that probability gives a yes or no.

\paragraph{What the data needed.} Only 1 person in 20 in the file had a
stroke. Trained as is, the model says ``no'' to almost everyone: 95\,\%
accuracy, but it finds 2 strokes out of 62. For a screening tool that is
useless. Telling the model that a missed stroke weighs as much as the
nineteen healthy people it is outnumbered by (balanced class weights)
fixes it.

\paragraph{What is kept.} The balanced model. It finds 50 of the 62
strokes, at the cost of sending about 300 healthy people out of 1\,216 to a
check-up. Age matters five times more than any other fact. Because only
about one positive in seven is a real stroke, the assistant presents a
positive answer as a reason to see a doctor, never as a diagnosis.

\subsection{Cirrhosis stage: support vector machine}

\paragraph{The question.} How advanced is this patient's liver disease,
from blood tests and physical signs, without a biopsy?

\paragraph{Why an SVM.} It draws the boundary between classes that leaves
the widest gap on each side, and with the radial kernel that boundary can
curve. Because it measures distances between patients, all the tests had to
be put on the same scale first, and its two settings ($C$ and $\gamma$)
were chosen by cross-validation.

\paragraph{What the model found.} Telling apart the four stages is not
possible from these data: stages 2 and 3 have the same typical values on
almost every test, and no setting of the model got above 45\,\% balanced
accuracy. What the tests do show is whether the patient has reached
cirrhosis (stage 4) or not.

\paragraph{What is kept.} The two-class model, cirrhosis or earlier, right
about three times in four. The four-stage model is saved next to it so the
API can still answer with a stage, with its limitation attached. The data
come from a clinical trial of the 1970s and 1980s, mostly women, so the
model is a screen, not a diagnosis.

\subsection{Avocado price: random forest}

\paragraph{The question.} What will an avocado cost in a given city, for
a given type (organic or conventional) and week?

\paragraph{Why a random forest.} A single decision tree memorises its
training data: grown without limits it is perfect on the rows it saw and
much worse on new ones. A forest grows hundreds of trees on random
samples and averages them, which keeps what the trees learn and cancels
what they memorise.

\paragraph{What the model found.} The forest of 200 trees explains 90\,\%
of the price with an average error of nine cents, and its out-of-bag
score agrees. The type of avocado is by far the most important fact
(organic costs about half a dollar more), followed by the quantities sold
and the season.

\paragraph{What is kept.} The forest, with one warning. Trained on
2015--2017 and asked about 2018, it explains only 41\,\% of the price,
because a forest cannot guess a year it has never seen. So it is used to
look up the expected price for a region, type and season, not to forecast
next year.

\subsection{Hepatitis C category: naive Bayes}

\paragraph{The question.} Does this blood panel look like a healthy
donor, hepatitis, fibrosis or cirrhosis?

\paragraph{Why naive Bayes.} It computes, for each category, how likely
the observed lab values are, and picks the most likely category. It works
well with few examples, which matters here: the file has 533 donors but
only 24, 21 and 30 sick patients.

\paragraph{What the data needed.} The Gaussian version assumes each lab
value follows a bell curve within each category. Liver enzymes do not:
most values are small and a few are huge. Taking the logarithm of every
lab value makes them bell-shaped and raises the balanced accuracy from
0.65 to 0.71.

\paragraph{What is kept.} The logarithmic model. It separates donors from
patients almost perfectly and recognises cirrhosis four times out of five,
but it confuses hepatitis and fibrosis with each other about half the
time, so the assistant reads either of those as ``early liver disease''.
Because the sick groups are so small, the numbers come from five-fold
cross-validation rather than a single test set.

\subsection{Body fat: multiple linear regression}

\paragraph{The question.} What is a person's body fat percentage, from
measurements taken with a tape and a scale?

\paragraph{Why linear regression.} The reference method (underwater
weighing) is expensive; a formula with a few readable coefficients is what
a trainer can actually use.

\paragraph{What the data needed.} The file has thirteen measurements, but
they mostly say the same thing: a bigger man has a bigger chest, waist,
hip and thigh, all at once. When predictors repeat each other, a
regression cannot decide how to share the effect between them and its
coefficients become meaningless, some even with the wrong sign. The
\code{Density} column was dropped because the target is computed from it.
Two impossible rows (a man 75 centimetres tall, a body fat of exactly 0)
were removed.

\paragraph{What is kept.} A model with three measurements: abdomen,
weight and wrist. It is as accurate as the model with all thirteen
(average error of 3.9 points of fat against 3.8, on a test set too small to
tell the difference) and its three coefficients can be read and trusted.
The limitation is the error itself: a man predicted at 20\,\% may be at 16
or 24, enough to place him in a broad category but not to follow a diet
week by week.
